\section{Datenerhebung: Experteninterviews}
\label{sec:04_interviews-data-collection}

Die Datenerhebung erfolgte mittels leitfadengestützter Experteninterviews in Kombination mit vergleichenden \Gls{usability}-Tests. Experteninterviews eignen sich insbesondere dann, wenn kontextgebundenes, erfahrungsbasiertes Wissen rekonstruiert werden soll, das für die Beantwortung der \Glspl{rq} zentral ist \cite{405:KvaleBrinkmann2015,420:GlaeserLaudel2010}. \citeauthor{420:GlaeserLaudel2010} fassen Experten dabei als Personen auf, die über spezifisches Wissen zu den zu rekonstruierenden Sachverhalten verfügen und dadurch in besonderer Weise zur Erklärung organisationaler Prozesse beitragen können \cite{420:GlaeserLaudel2010}.

Die Interviewpartner gelten in diesem Sinne als Experten, da sie entweder regelmäßig mit dem bestehenden Wiki arbeiten, Verantwortung für Inhalte oder Strukturen tragen oder als Vertreter relevanter Zielgruppen fungieren. Durch die Auswahl verschiedener Rollen (\acrshort{zB} operative Projektmitarbeitende, koordinierende Funktionen, Führungskräfte) werden unterschiedliche Nutzungsperspektiven auf die Plattform einbezogen und damit zentrale Akteursgruppen der Fallstudie abgedeckt \cite{404:Yin2018}. 

\newpage

Die Kombination von Experteninterviews und \Gls{usability}-Tests erlaubt es, sowohl subjektive Einschätzungen als auch tatsächliches Nutzungsverhalten systematisch zu erfassen. Während die \Gls{usability}-Tests konkrete Interaktionsabläufe im alten und neuen Wiki sichtbar machen, dienen die anschließenden Interviews dazu, diese Erfahrungen zu kontextualisieren, zu interpretieren und mit organisationsbezogenen Erwartungen sowie Rahmenbedingungen zu verknüpfen \cite{405:KvaleBrinkmann2015}.

Dieses Vorgehen folgt einem Mixed-Methods-Forschungsansatz, bei dem qualitative und quantitative Erhebungsformen gezielt kombiniert werden, um unterschiedliche Erkenntnisperspektiven auf den Untersuchungsgegenstand zu ermöglichen. Abbildung~\ref{fig:mixed-methods-triangulation} verdeutlicht die Logik dieser methodischen Triangulation. Quantitative Usability-Metriken liefern strukturierte Hinweise auf Effizienz, Effektivität und subjektive Zufriedenheit, während qualitative Daten aus Interviews, Beobachtungen und Protokollen ein vertieftes Verständnis der Nutzungserfahrungen und Bedeutungszuschreibungen der Beteiligten ermöglichen. Die Zusammenführung dieser Daten erfolgt auf interpretativer Ebene und dient der wechselseitigen Einordnung, Plausibilisierung und Kontextualisierung der Ergebnisse.

\image{H}{1}{g-12_mixed-methods-triangulation}{
Mixed-Methods-Triangulation der Datenerhebung durch qualitative und quantitative Erhebungsformen sowie deren interpretative Zusammenführung
}

Die standardisierte Feedbackumfrage ergänzt dieses Vorgehen um vergleichbare Einschätzungen zu Nutzbarkeit, Verständlichkeit und Akzeptanz. Die erhobenen quantitativen Daten dienen dabei nicht der statistischen Hypothesenprüfung, sondern unterstützen im Sinne eines überwiegend qualitativen Erkenntnisinteresses die Strukturierung und Veranschaulichung der qualitativen Befunde \cite{200:ScribbrInductiveDeductive2023}. Insgesamt ermöglicht der Mixed-Methods-Ansatz eine differenzierte Betrachtung des Wiki-Auftritts, bei der beobachtbares Nutzungsverhalten, subjektive Bewertungen und organisationale Deutungsmuster systematisch miteinander in Beziehung gesetzt werden.


\subsection{Auswahl der Interviewpartner}
\label{subsec:04-02-01_sampling-and-selection}
\subsection{Leitfadenentwicklung}
\label{subsec:04-02-02_interview-guide-development}
\subsection{Durchführung und Dokumentation}
\label{subsec:04-02-03_execution-and-documentation}